% simplicial_complex_embedders_v3.tex
%
% A rewrite of v1 that presents intent first and algorithm second.
% The algorithm steps are introduced as the natural implementation of
% a geometric strategy, not as a sequence of operations to be explained
% after the fact.
%
% v3 changes (vs v2): the Algorithm-2 worked-example tables are no longer
% hand-coded.  They are generated from EncDec.py by
% generate_simplex_example_tables.py and \input from DOCS/generated/.
% To change the example multiset, edit M in that script and re-run
% (make tables).  Steps C and D are also now presented as two instances
% of the SAME Abel-summation operation, in matching table formats.
% ═════════════════════════════════════════════════════════════════════════════
\section{What We Are Trying to Do}
\label{sec:goals}
% ═════════════════════════════════════════════════════════════════════════════

The input is an unordered multiset $M = \{\mathbf{v}_0,\ldots,\mathbf{v}_{n-1}\}$
of $n$ vectors in $\R^k$.  Concretely we represent $M$ as a $k\times n$ matrix
whose columns are the vectors $\mathbf{v}_j$.  Column order is arbitrary.

We want a map $f$ from such multisets to some $\R^N$ that is:
\begin{itemize}[nosep]
  \item \textbf{Multiset-invariant} — the output does not depend on how the
        columns are ordered.
  \item \textbf{Injective} — distinct multisets give distinct outputs ($f$ is
        an \emph{embedder}).
  \item \textbf{Continuous} ($C^0$) — small changes in the vectors produce
        small changes in $f(M)$.
  \item \textbf{Positively homogeneous of degree 1} — $f(\lambda M) =
        \lambda f(M)$ for all $\lambda > 0$.
\end{itemize}

Multiset invariance and continuity are in tension: any map that achieves
invariance by sorting will have kinks wherever the sorted order changes.
The strategy below resolves this tension elegantly.

% ═════════════════════════════════════════════════════════════════════════════
\section{The Strategy}
\label{sec:strategy}
% ═════════════════════════════════════════════════════════════════════════════

\subsection*{Step A — Write $M$ in a natural basis}

Label each of the $nk$ matrix entries by $\eji{j}{i}$ (column $j$, row $i$).
Let $X[j][i]$ denote the entry in column $j$, row $i$.  Then
\[
  M \;=\; \sum_{j,i} X[j][i]\;\eji{j}{i},
\]
There is no harm in viewing the $\eji{j}{i}$ as totally abstract/opaque basis vectors with no physical intepretation, but sometimes it can be helpful to view $\eji{j}{i}$ concretely as the $k\times n$ matrix that is 1 at $(i,j)$ and 0
elsewhere.  This is the expansion of $M$ in the \emph{standard Eji basis}.

\subsection*{Step B — Peel off the fully symmetric parts}

A \emph{fully symmetric} basis element at row $i$ is $\sum_j \eji{j}{i}$: it
places the same coefficient on every vector at coordinate $i$, so it is
unaffected by any permutation of the columns.  The smallest coefficient that
any vector has at row $i$ is $m_i = \min_j X[j][i]$.  Factoring:
\[
  M \;=\; \underbrace{\sum_i m_i \!\left(\sum_j \eji{j}{i}\right)}_{\text{symmetric part}}
    \;+\; \underbrace{\text{balance}}_{\text{all entries} \;\ge\; 0}.
\]
The symmetric part is already permutation-invariant; its $k$ scalar
coefficients $m_0,\ldots,m_{k-1}$ can be stored directly as anchor values.

The balance $M - \sum_i m_i (\sum_j \eji{j}{i})$ has all non-negative
entries and contains the non-trivial, non-symmetric information in $M$.

\subsection*{Step C — Re-express the balance as a positive combination of
             nested basis elements}

Within row $i$, order the $n$ vectors descending by their (shifted) value.
Let $\sigma_i(0),\ldots,\sigma_i(n-1)$ be this ordering.  The gaps
$\delta^{(i)}_r = X[\sigma_i(r)][i] - X[\sigma_i(r+1)][i] \ge 0$ are the
drops between consecutive ranks.

By Abel summation (the discrete analogue of integration by parts; we also call
this a \emph{barycentric subdivision}, or BSD), the balance
at row $i$ becomes:
\[
  \text{balance at row }i
    \;=\; \sum_{r=0}^{n-2} \delta^{(i)}_r \cdot V^{(i)}_r,
\]
where the \emph{prefix basis element} $V^{(i)}_r =
\bigl\{\eji{\sigma_i(0)}{i},\ldots, \eji{\sigma_i(r)}{i}\bigr\}$
is the set of the top-$(r{+}1)$ vectors by row-$i$ value.

This is another change of basis: from individual Ejis with signed coefficients
to nested prefix sets with non-negative gap coefficients.  The structure of
which vectors appear in $V^{(i)}_r$ — which $j$'s dominate row $i$ — is the
non-symmetric information not yet discarded.

\subsection*{Step D — Merge, re-sort, and apply a second Abel summation}

Collect all $m = k(n-1)$ gap--prefix-element pairs from all rows into one list
and sort descending by gap: $\delta_{(0)} \ge \cdots \ge \delta_{(m-1)}$ with
corresponding basis elements $V_{(0)},\ldots,V_{(m-1)}$.

Apply a second Abel summation to obtain \emph{cumulative} basis elements and
weighted-difference coefficients:
\[
  \alpha_s \;=\; (s+1)\bigl(\delta_{(s)} - \delta_{(s+1)}\bigr),
  \qquad \delta_{(m)} := 0,
\]
\[
  L_s \;=\; V_{(0)} + V_{(1)} + \cdots + V_{(s)},
\]
where $L_s$ is the \emph{Eji\_LinComb} (also called a \emph{snapshot}; see
\S\ref{sec:bridge}): the $k\times n$ matrix whose $(i,j)$
entry counts how many of the first $s+1$ prefix elements contain $\eji{j}{i}$.
Its \emph{index} $s+1$ records how many prefix elements have been accumulated.

The mathematical identity $\sum_s \delta_{(s)} V_{(s)} = \sum_s \alpha_s
\cdot (L_s/(s+1))$ confirms this preserves all information.  The normalised
basis element $L_s/(s+1)$ is the \emph{average} of the first $s+1$ prefix
elements.  The $(s+1)$ weight in $\alpha_s$ distributes contributions more
evenly and ensures that $\sum_s \alpha_s = \sum_s \delta_{(s)}$ (the total
weight is preserved).

Crucially, $L_s$ captures \emph{which prefix elements came first} in the
global sort: two different orderings of the same set of prefix elements produce
different $L_s$ sequences.

\subsection*{Step E — Achieve permutation invariance by canonicalising}

The $j$-indices in $L_s$ still reflect the arbitrary labelling of columns.
The \emph{canonical form} $\canon{L_s}$ is obtained by sorting the columns of
$L_s$'s count matrix lexicographically — this removes the column labelling
while retaining the multiset structure.  Two inputs related by a column
permutation produce identical $\canon{L_s}$ sequences and hence identical
outputs.
%
% Forward hook (added for the combined document; no original sentence reworded):
% the claim that this per-matrix LOCAL canonicalisation is injective-safe is
% exactly what \S\ref{sec:bridge} states and the part after it proves.
That each $L_s$ is canonicalised \emph{on its own}---a \emph{local} column
sort that discards how the columns of one $L_s$ correspond to those of
another---is what makes this step cheap and continuous; that it nonetheless
cannot make two genuinely different multisets collide is the one substantive
ingredient of injectivity, stated and proved in \S\ref{sec:bridge}
(\emph{Injectivity}).

\subsection*{Step F — Encode in $\R^{2m+1}$ and assemble}

Map each canonical $L_s$ to a pseudorandom point $p_s = \hash{\canon{L_s}}
\in [0,1]^{2m+1}$ via MD5 hashing of its count matrix and index.  The main
body of the output is $\sum_s \alpha_s p_s$.

\begin{tcolorbox}[keybox]
\textbf{Why $\R^{2m+1}$ is the right dimension.}
All $\alpha_s \ge 0$, so $\sum_s \alpha_s p_s$ is a point in the
\emph{positive cone} generated by $\{p_0,\ldots,p_{m-1}\}$, a cone of
dimension $m$.  Two $m$-dimensional positive cones in $\R^{2m+1}$ are
\emph{generically disjoint} by dimension counting: $m + m = 2m < 2m+1$.
Therefore a single point in $\R^{2m+1}$ encodes both:
(a) which cone it lies in (i.e.\ which canonical $\{L_s\}$, i.e.\ which basis
elements), and (b) the coefficients $\{\alpha_s\}$ within that cone.

This is also why the algorithm is $C^0$: when two values in the same row are
equal, the gap at that position is $0$ and so $\alpha_s = 0$.  The point
$\sum \alpha_s p_s$ then lies on a shared \emph{face} of two adjacent cones
rather than in either interior.  That shared face is the geometric expression
of the continuity: the change in which cone we are in happens exactly at the
moment when the coefficient of the transitioning basis element is zero, so the
output is unaffected.
\end{tcolorbox}

\noindent
The full output is:
\[
  f(M) \;=\; \Bigl[\;\underbrace{m_0,\ldots,m_{k-1}}_{k\;\text{row minima}},\;
                     \underbrace{\textstyle\sum_{s=0}^{m-1}\alpha_s p_s}_{2m+1\;\text{reals}}
             \;\Bigr]
  \;\in\; \R^{k + 2m + 1}.
\]
With $m = k(n-1)$: total size $= k + 2k(n-1) + 1 = 2nk+1-k$.

% ═════════════════════════════════════════════════════════════════════════════
\section{Two Algorithms, One Strategy}
\label{sec:two-algs}
% ═════════════════════════════════════════════════════════════════════════════

Both algorithms follow Steps A--F.  They differ only in \textbf{Step B}: how
many symmetric parts are peeled off, and what ``fully symmetric'' means.

\paragraph{Algorithm 2 (per-row).}
Factors out the $k$ per-row symmetric elements $\sum_j \eji{j}{i}$, one per
coordinate.  Stores $k$ row minima as anchors.  Steps C--F operate on
$m = k(n-1)$ gaps.
Implemented as \texttt{Embedder.embed\_generic} in \algtwoapy{} (direct) and
\algtwobpy{} (EncDec-backed); the \algtwopy{} wrapper selects between them.
See also \texttt{simplex\_2\_preprocess\_steps} in \encdecpy{} and
Appendix~\ref{app:code}.

\paragraph{Algorithm 1 (global).}
Factors out only the \emph{single} fully-symmetric element $\sum_{i,j}
\eji{j}{i}$ (treating every entry identically), with coefficient = global
minimum.  The global maximum is also recorded (see note below).  Steps C--F
then operate on a single globally-sorted list of $m = nk-1$ gaps.
Implemented as \texttt{Embedder.embed\_generic} in \algonepy{} and as
\texttt{simplex\_1\_preprocess\_steps} in \encdecpy{}; see
Appendix~\ref{app:code}.

The single global sort of Algorithm~1 means that prefix elements can contain
Ejis from \emph{any} row, whereas in Algorithm~2 every prefix element contains
Ejis from exactly one row.  This is why Algorithm~2 is conceptually cleaner:
it respects the coordinate structure of the vectors.

\paragraph{Dimension count.}
\begin{center}
\renewcommand{\arraystretch}{1.3}
\begin{tabular}{lcc}
  \toprule
  & \textbf{\textcolor{alg1col}{Algorithm 1}} & \textbf{\textcolor{alg2col}{Algorithm 2}} \\
  \midrule
  Symmetric parts extracted & 1 (global) & $k$ (per-row) \\
  Gaps $m$ & $nk-1$ & $k(n-1)=nk-k$ \\
  Anchor values stored & 2 (max \& min) & $k$ (row minima) \\
  Hash embedding dimension $2m+1$ & $2nk-1$ & $2nk-2k+1$ \\
  \textbf{Total output size} & $\mathbf{2nk+1}$ & $\mathbf{2nk+1-k}$ \\
  \bottomrule
\end{tabular}
\end{center}

Algorithm~2 saves $k$ output dimensions because per-row symmetric extraction
produces $k-1$ fewer gaps than global extraction (since each row contributes
$n-1$ gaps, totalling $k(n-1) = nk-k$, versus the global $nk-1$ gaps from
Algorithm~1), and the $(s+1)$ weighting scheme in Step~D ensures that the
hash dimension savings exceed the extra anchor overhead.

\begin{tcolorbox}[notebox]
\textbf{Note on Algorithm 1's global maximum.}  In a perfect implementation,
only the global minimum is needed as an anchor (it fixes the absolute scale;
all relative gaps are encoded in the hash sum).  The global maximum is
technically redundant.  It is included in Algorithm~1 for implementation
simplicity; a TODO comment in the source code acknowledges this.
\end{tcolorbox}

\begin{tcolorbox}[keybox]
\textbf{Source code.}

\smallskip
\noindent\textit{Algorithm~2 — two interchangeable implementations (Steps~A--F each):}
\begin{itemize}[nosep,leftmargin=1.5em,topsep=2pt]
  \item \algtwoapy{} — direct implementation; \texttt{Embedder.embed\_generic}.
    This is the original file on which the document is based.
  \item \algtwobpy{} — EncDec-backed; \texttt{Embedder.embed\_generic}.
    Delegates Steps~A--E to \encdecpy{} and adds Step~F directly.
    Default config produces bit-identical output to \algtwoapy{}.
  \item \algtwopy{} — wrapper that selects between the above.
    All other code imports \texttt{Embedder} from this file.
\end{itemize}

\smallskip
\noindent\textit{Algorithm~1:}
\begin{itemize}[nosep,leftmargin=1.5em,topsep=2pt]
  \item \algonepy{} — direct implementation; \texttt{Embedder.embed\_generic}; Steps~A--F.
\end{itemize}

\smallskip
\noindent\textit{Shared infrastructure:}
\begin{itemize}[nosep,leftmargin=1.5em,topsep=2pt]
  \item \encdecpy{} — \texttt{simplex\_2\_preprocess\_steps} and
    \texttt{simplex\_1\_preprocess\_steps} cover Steps~A--E.
    Docstrings use this document's $\delta_{(s)}$, $V_{(s)}$, $\Delta_{(s)}$,
    $L_{(s)}$, $\alpha_s$, $\hat{L}_{(s)}$ notation.
  \item \playpy{} — command-line driver for \encdecpy{}.
\end{itemize}

\smallskip
\noindent A parity test suite
(\texttt{test\_simplex2\_ab\_parity.py}) verifies that \algtwoapy{} and
\algtwobpy{} produce identical embeddings under the default configuration.
See Appendix~\ref{app:code} for a full notation--code correspondence table.
\end{tcolorbox}

% ═════════════════════════════════════════════════════════════════════════════
\section{Worked Example: Algorithm 2, $n=4$, $k=3$}
\label{sec:example-alg2}
% ═════════════════════════════════════════════════════════════════════════════

All tables and numeric displays in this section are generated directly from
\encdecpy{} by \texttt{DOCS/generate\_simplex\_example\_tables.py} and
\texttt{\textbackslash input} from \texttt{DOCS/generated/}.  To re-run the
example with a different multiset, edit $M$ in that script and rebuild
(\texttt{make tables}).

% AUTO-GENERATED by generate_simplex_example_tables.py — do not edit.
% Regenerate with: make tables   (or run the script directly).
\[
  M \;=\; \left\{\;
    \begin{pmatrix}4\\2\\3\end{pmatrix};\;
    \begin{pmatrix}-3\\-5\\1\end{pmatrix};\;
    \begin{pmatrix}8\\9\\2\end{pmatrix};\;
    \begin{pmatrix}2\\-3\\-7\end{pmatrix}
  \;\right\}
  \quad (n=4 \text{ vectors in } \R^3)
\]


\subsection*{Step B: Extract row minima}

Row minima: % AUTO-GENERATED by generate_simplex_example_tables.py — do not edit.
% Regenerate with: make tables   (or run the script directly).
$m_0 = -3$ (row~0), $m_1 = 2$ (row~1), $m_2 = 1$ (row~2)%
.
These three scalars are stored as anchors.  After subtracting:
% AUTO-GENERATED by generate_simplex_example_tables.py — do not edit.
% Regenerate with: make tables   (or run the script directly).
\[
  \text{balance} = \left\{\;
    \begin{pmatrix}7\\0\\2\end{pmatrix};\;
    \begin{pmatrix}0\\3\\0\end{pmatrix};\;
    \begin{pmatrix}11\\7\\1\end{pmatrix};\;
    \begin{pmatrix}5\\5\\1\end{pmatrix}
  \;\right\}
\]

(All non-negative.)  The debug check in the source code verifies that
$\sum_s \alpha_s \cdot (L_s/(s+1)) + \text{broadcast}(m_i)$ recovers
the original $M$ exactly.

\subsection*{Step C: First Abel summation (once per row)}

\noindent
Each of the $k=3$ coordinate rows is now fed through the Abel-summation
operation — the \emph{same} operation that Step~D will apply once more to the
merged results, and the same operation that \encdecpy{} implements as
\texttt{barycentric\_subdivide}.  Its input is a list of (coefficient, basis
element) pairs; its output is a list of (gap, cumulative prefix) pairs plus
one \emph{offset} term.  The three tables below (one per row) deliberately use
the same layout as the large Step-D table: the left column pair is the input,
the right column pair is the output.

Each row's input is its raw values $x^{(i)}_{(r)}$ paired with single Ejis;
sorting descending and differencing gives the gaps
$\delta^{(i)}_r = x^{(i)}_{(r)} - x^{(i)}_{(r+1)}$ and cumulative prefixes
$V^{(i)}_r$.  Note that the gaps are unchanged by the Step-B subtraction, and
the offset term of each table — final value $\times$ all-ones row — is
\emph{precisely} the row-minimum anchor $m_i$ of Step~B.  (In \encdecpy{},
Step~B is not a separate pass: the anchors are simply the offset outputs of
these three calls, requested via \texttt{return\_offset\_separately=True}.)

% AUTO-GENERATED by generate_simplex_example_tables.py — do not edit.
% Regenerate with: make tables   (or run the script directly).
\begin{center}
\small
\setlength{\tabcolsep}{4pt}
\renewcommand{\arraystretch}{1.35}
\begin{tabular}{c|cl|cl}
\toprule
\multicolumn{5}{l}{\textbf{Row $i=0$}} \\
\midrule
  & & & $x^{(0)}_{(r)}-x^{(0)}_{(r+1)}$ & $V^{(0)}_{r-1}+\eji{\sigma_0(r)}{0}$ \\
\midrule
$r$ & $x^{(0)}_{(r)}$ & $\eji{\sigma_0(r)}{0}$ & $\delta^{(0)}_r$ & $V^{(0)}_r$ \\
\midrule
0 & $8$ & $\eji{2}{0}$ & 4
  & $\eji{2}{0}$
  \\[2pt]
1 & $4$ & $\eji{0}{0}$ & 2
  & $\eji{0}{0}+\eji{2}{0}$
  \\[2pt]
2 & $2$ & $\eji{3}{0}$ & 5
  & $\eji{0}{0}+\eji{2}{0}+\eji{3}{0}$
  \\[2pt]
\midrule
\multicolumn{5}{l}{offset $\;=\;x^{(0)}_{(3)}\cdot\textstyle\sum_j\eji{j}{0}\;=\;-3\,(\eji{0}{0}+\eji{1}{0}+\eji{2}{0}+\eji{3}{0})$} \\
\bottomrule
\end{tabular}
\end{center}

\begin{center}
\small
\setlength{\tabcolsep}{4pt}
\renewcommand{\arraystretch}{1.35}
\begin{tabular}{c|cl|cl}
\toprule
\multicolumn{5}{l}{\textbf{Row $i=1$}} \\
\midrule
  & & & $x^{(1)}_{(r)}-x^{(1)}_{(r+1)}$ & $V^{(1)}_{r-1}+\eji{\sigma_1(r)}{1}$ \\
\midrule
$r$ & $x^{(1)}_{(r)}$ & $\eji{\sigma_1(r)}{1}$ & $\delta^{(1)}_r$ & $V^{(1)}_r$ \\
\midrule
0 & $9$ & $\eji{2}{1}$ & 2
  & $\eji{2}{1}$
  \\[2pt]
1 & $7$ & $\eji{3}{1}$ & 2
  & $\eji{2}{1}+\eji{3}{1}$
  \\[2pt]
2 & $5$ & $\eji{1}{1}$ & 3
  & $\eji{1}{1}+\eji{2}{1}+\eji{3}{1}$
  \\[2pt]
\midrule
\multicolumn{5}{l}{offset $\;=\;x^{(1)}_{(3)}\cdot\textstyle\sum_j\eji{j}{1}\;=\;2\,(\eji{0}{1}+\eji{1}{1}+\eji{2}{1}+\eji{3}{1})$} \\
\bottomrule
\end{tabular}
\end{center}

\begin{center}
\small
\setlength{\tabcolsep}{4pt}
\renewcommand{\arraystretch}{1.35}
\begin{tabular}{c|cl|cl}
\toprule
\multicolumn{5}{l}{\textbf{Row $i=2$}} \\
\midrule
  & & & $x^{(2)}_{(r)}-x^{(2)}_{(r+1)}$ & $V^{(2)}_{r-1}+\eji{\sigma_2(r)}{2}$ \\
\midrule
$r$ & $x^{(2)}_{(r)}$ & $\eji{\sigma_2(r)}{2}$ & $\delta^{(2)}_r$ & $V^{(2)}_r$ \\
\midrule
0 & $3$ & $\eji{0}{2}$ & 1
  & $\eji{0}{2}$
  \\[2pt]
1 & $2$ & $\eji{2}{2}$ & 0
  & $\eji{0}{2}+\eji{2}{2}$
  \\[2pt]
2 & $2$ & $\eji{3}{2}$ & 1
  & $\eji{0}{2}+\eji{2}{2}+\eji{3}{2}$
  \\[2pt]
\midrule
\multicolumn{5}{l}{offset $\;=\;x^{(2)}_{(3)}\cdot\textstyle\sum_j\eji{j}{2}\;=\;1\,(\eji{0}{2}+\eji{1}{2}+\eji{2}{2}+\eji{3}{2})$} \\
\bottomrule
\end{tabular}
\end{center}


\noindent
Normalisation (\texttt{preserve\_scale}) is switched \emph{off} at this layer
in the standard configuration, so no $\alpha$/$\hat{V}$ column pair appears;
the \algtwobpy{} implementation also permits it to be switched on here, which
would add one, exactly as in Step~D's table.

\subsection*{Step D: Merge, re-sort, and the same Abel summation again}

The nine $(\delta^{(i)}_r, V^{(i)}_r)$ output pairs of Step~C now become the
\emph{input} pairs $(\delta_{(s)}, V_{(s)})$ of a second pass through the very
same operation.  Merge all $9 = k(n{-}1)$ pairs and sort by gap descending:
% AUTO-GENERATED by generate_simplex_example_tables.py — do not edit.
% Regenerate with: make tables   (or run the script directly).
\[
  5,\;4,\;3,\;2,\;2,\;2,\;1,\;1,\;0
\]

The table below has the same structure as Step~C's tables — input pair on the
left, output pair(s) on the right — but this time with normalisation switched
\emph{on}, which appends the $(\alpha_s, \hat{L}_{(s)})$ column pair to the
unnormalised $(\Delta_{(s)}, L_{(s)})$ one.  Terms within each $L_s$
expression are grouped by coordinate row $i$.

% AUTO-GENERATED by generate_simplex_example_tables.py — do not edit.
% Regenerate with: make tables   (or run the script directly).
\begin{center}
\small
\setlength{\tabcolsep}{4pt}
\renewcommand{\arraystretch}{1.35}
\begin{tabular}{c|cp{3cm}|cp{3cm}|cp{3cm}}
\toprule
    &
    &
    & $\delta_{(s)}-\delta_{(s+1)}$
    & $L_{(s-1)}+V_{(s)}$
    & $(s+1)\Delta_{(s)}$
    & $L_{(s)}/(s+1)$\\
\midrule
$s$ & $\delta_{(s)}$
    & $V_{(s)}$
    & $\Delta_{(s)}$
    & $L_{(s)}$
    & $\alpha_s$
    & ${\hat L}_{(s)}$ \\
\midrule
0 & $ 8 $ & $\eji{0}{2}+\eji{1}{2}+\eji{2}{2}$ & $ 1 $
  & $\eji{0}{2}+\eji{1}{2}+\eji{2}{2}$
    & $ 1 $
  & $\eji{0}{2}+\eji{1}{2}+\eji{2}{2}$
   \\[3pt]
1 & $ 7 $ & $\eji{2}{1}$ & $ 2 $
  & $\eji{2}{1}$\newline
    ${}+\eji{0}{2}+\eji{1}{2}+\eji{2}{2}$
    & $ 4 $
  & $\tfrac{1}{2} (\eji{2}{1}$\newline
    ${}+\eji{0}{2}+\eji{1}{2}+\eji{2}{2})$
   \\[3pt]
\rowcolor{gray!15}\cellcolor{white}2 & \textcolor{gray}{$ 5 $} & \textcolor{gray}{$\eji{0}{0}+\eji{2}{0}+\eji{3}{0}$} & \textcolor{gray}{$ 0 $}
  & \textcolor{gray}{$\eji{0}{0}+\eji{2}{0}+\eji{3}{0}$\newline
    ${}+\eji{2}{1}$\newline
    ${}+\eji{0}{2}+\eji{1}{2}+\eji{2}{2}$}
    & \textcolor{gray}{$ 0 $}
  & \textcolor{gray}{$\tfrac{1}{3} (\eji{0}{0}+\eji{2}{0}+\eji{3}{0}$\newline
    ${}+\eji{2}{1}$\newline
    ${}+\eji{0}{2}+\eji{1}{2}+\eji{2}{2})$}
   \\[3pt]
3 & \cellcolor{gray!15}\textcolor{gray}{$ 5 $} & \cellcolor{gray!15}\textcolor{gray}{$\eji{0}{1}+\eji{2}{1}$} & $ 1 $
  & $\eji{0}{0}+\eji{2}{0}+\eji{3}{0}$\newline
    ${}+\eji{0}{1}+2\eji{2}{1}$\newline
    ${}+\eji{0}{2}+\eji{1}{2}+\eji{2}{2}$
    & $ 4 $
  & $\tfrac{1}{4} (\eji{0}{0}+\eji{2}{0}+\eji{3}{0}$\newline
    ${}+\eji{0}{1}+2\eji{2}{1}$\newline
    ${}+\eji{0}{2}+\eji{1}{2}+\eji{2}{2})$
   \\[3pt]
4 & $ 4 $ & $\eji{2}{0}$ & $ 2 $
  & $\eji{0}{0}+2\eji{2}{0}+\eji{3}{0}$\newline
    ${}+\eji{0}{1}+2\eji{2}{1}$\newline
    ${}+\eji{0}{2}+\eji{1}{2}+\eji{2}{2}$
    & $ 10 $
  & $\tfrac{1}{5} (\eji{0}{0}+2\eji{2}{0}+\eji{3}{0}$\newline
    ${}+\eji{0}{1}+2\eji{2}{1}$\newline
    ${}+\eji{0}{2}+\eji{1}{2}+\eji{2}{2})$
   \\[3pt]
\rowcolor{gray!15}\cellcolor{white}5 & \textcolor{gray}{$ 2 $} & \textcolor{gray}{$\eji{0}{0}+\eji{2}{0}$} & \textcolor{gray}{$ 0 $}
  & \textcolor{gray}{$2\eji{0}{0}+3\eji{2}{0}+\eji{3}{0}$\newline
    ${}+\eji{0}{1}+2\eji{2}{1}$\newline
    ${}+\eji{0}{2}+\eji{1}{2}+\eji{2}{2}$}
    & \textcolor{gray}{$ 0 $}
  & \textcolor{gray}{$\tfrac{1}{6} (2\eji{0}{0}+3\eji{2}{0}+\eji{3}{0}$\newline
    ${}+\eji{0}{1}+2\eji{2}{1}$\newline
    ${}+\eji{0}{2}+\eji{1}{2}+\eji{2}{2})$}
   \\[3pt]
6 & \cellcolor{gray!15}\textcolor{gray}{$ 2 $} & \cellcolor{gray!15}\textcolor{gray}{$\eji{0}{1}+\eji{2}{1}+\eji{3}{1}$} & $ 1 $
  & $2\eji{0}{0}+3\eji{2}{0}+\eji{3}{0}$\newline
    ${}+2\eji{0}{1}+3\eji{2}{1}+\eji{3}{1}$\newline
    ${}+\eji{0}{2}+\eji{1}{2}+\eji{2}{2}$
    & $ 7 $
  & $\tfrac{1}{7} (2\eji{0}{0}+3\eji{2}{0}+\eji{3}{0}$\newline
    ${}+2\eji{0}{1}+3\eji{2}{1}+\eji{3}{1}$\newline
    ${}+\eji{0}{2}+\eji{1}{2}+\eji{2}{2})$
   \\[3pt]
\rowcolor{gray!15}\cellcolor{white}7 & \textcolor{gray}{$ 1 $} & \textcolor{gray}{$\eji{0}{2}$} & \textcolor{gray}{$ 0 $}
  & \textcolor{gray}{$2\eji{0}{0}+3\eji{2}{0}+\eji{3}{0}$\newline
    ${}+2\eji{0}{1}+3\eji{2}{1}+\eji{3}{1}$\newline
    ${}+2\eji{0}{2}+\eji{1}{2}+\eji{2}{2}$}
    & \textcolor{gray}{$ 0 $}
  & \textcolor{gray}{$\tfrac{1}{8} (2\eji{0}{0}+3\eji{2}{0}+\eji{3}{0}$\newline
    ${}+2\eji{0}{1}+3\eji{2}{1}+\eji{3}{1}$\newline
    ${}+2\eji{0}{2}+\eji{1}{2}+\eji{2}{2})$}
   \\[3pt]
8 & \cellcolor{gray!15}\textcolor{gray}{$ 1 $} & \cellcolor{gray!15}\textcolor{gray}{$\eji{0}{2}+\eji{2}{2}$} & $ 1 $
  & $2\eji{0}{0}+3\eji{2}{0}+\eji{3}{0}$\newline
    ${}+2\eji{0}{1}+3\eji{2}{1}+\eji{3}{1}$\newline
    ${}+3\eji{0}{2}+\eji{1}{2}+2\eji{2}{2}$
    & $ 9 $
  & $\tfrac{1}{9} (2\eji{0}{0}+3\eji{2}{0}+\eji{3}{0}$\newline
    ${}+2\eji{0}{1}+3\eji{2}{1}+\eji{3}{1}$\newline
    ${}+3\eji{0}{2}+\eji{1}{2}+2\eji{2}{2})$
   \\[3pt]
\midrule
&\multicolumn{2}{p{3.5cm}|}{$\sum_s \delta_{(s)}=35$}
  & \multicolumn{2}{p{3.5cm}|}{$\sum_s \Delta_{(s)}=8\ne 35$}
  & \multicolumn{2}{p{3.5cm}}{$\sum_s \alpha_{(s)}=35$}
  \\
     \midrule
& \multicolumn{2}{p{3.5cm}|}{
   $\sum_s \delta_{(s)} V_{(s)}=$}
  &\multicolumn{2}{p{3.5cm}|}{$
  \sum_s \Delta_{(s)} L_{(s)}=$}
  &\multicolumn{2}{p{3.5cm}}{$
  \sum_s \alpha_{(s)} {\hat L}_{(s)}=$}
   \\
&
&
   $7\eji{0}{0}+11\eji{2}{0}+5\eji{3}{0}$\newline
    ${}+7\eji{0}{1}+14\eji{2}{1}+2\eji{3}{1}$\newline
    ${}+10\eji{0}{2}+8\eji{1}{2}+9\eji{2}{2}$
  &
  &
  $7\eji{0}{0}+11\eji{2}{0}+5\eji{3}{0}$\newline
    ${}+7\eji{0}{1}+14\eji{2}{1}+2\eji{3}{1}$\newline
    ${}+10\eji{0}{2}+8\eji{1}{2}+9\eji{2}{2}$
  &
  & $7\eji{0}{0}+11\eji{2}{0}+5\eji{3}{0}$\newline
    ${}+7\eji{0}{1}+14\eji{2}{1}+2\eji{3}{1}$\newline
    ${}+10\eji{0}{2}+8\eji{1}{2}+9\eji{2}{2}$
  \\
\bottomrule
\end{tabular}
\end{center}


%\noindent
%Note how $L_s$ grows by exactly $V_{(s)}$ at each step: newly introduced Ejis
%appear with coefficient~1; Ejis already present have their count incremented.
%$L_s$~with index~$s{+}1$ is thus the count of how many of the first $s{+}1$
%prefix elements contain each $\eji{j}{i}$.

\begin{tcolorbox}[notebox]
\textbf{Reading off an $\alpha_s$.}  Each $\alpha_s$ is the product of the
position weight $(s{+}1)$ and the gap difference
$\Delta_{(s)} = \delta_{(s)}-\delta_{(s+1)}$, either of which can exceed~1.
For example:
% AUTO-GENERATED by generate_simplex_example_tables.py — do not edit.
% Regenerate with: make tables   (or run the script directly).
\[
  \alpha_{1} \;=\; (1{+}1)\bigl(\delta_{(1)}-\delta_{(2)}\bigr) \;=\; 2\times(7-5) \;=\; 4.
\]

(The example multiset used in earlier versions of this document had, by
coincidence, every non-zero gap difference equal to 1, making the non-zero
$\alpha_s$ values misleadingly equal to $s{+}1$.  The present multiset was
chosen to avoid that.)
\end{tcolorbox}

\begin{tcolorbox}[notebox]
\textbf{Verification: column sums and three-way equality.}
The integer columns $\delta_{(s)}$ and $\alpha_s$ have the same total:
% AUTO-GENERATED by generate_simplex_example_tables.py — do not edit.
% Regenerate with: make tables   (or run the script directly).
\[
  \sum_{s=0}^{8}\delta_{(s)} \;=\; 5+4+3+2+2+2+1+1+0 \;=\; 20
  \;=\; 1+2+3+0+0+6+0+8+0 \;=\; \sum_{s=0}^{8}\alpha_s.
\]

The three weighted Eji sums are all equal and recover the
\texttt{balance} matrix of Step~B (the $L_{(s)}$ and ${\hat L}_{(s)}$ totals-row
entries show the value directly; $\sum_s\delta_{(s)}V_{(s)}$ is identical):
% AUTO-GENERATED by generate_simplex_example_tables.py — do not edit.
% Regenerate with: make tables   (or run the script directly).
\begin{align*}
  \sum_s \delta_{(s)}\,V_{(s)}
  \;=\; \sum_s \Delta_{(s)}\,L_{(s)}
  \;=\; \sum_s \alpha_s\,{\hat L}_{(s)}
  &= 7\eji{0}{0}+11\eji{2}{0}+5\eji{3}{0} \\
  &\quad{}+3\eji{1}{1}+7\eji{2}{1}+5\eji{3}{1} \\
  &\quad{}+2\eji{0}{2}+\eji{2}{2}+\eji{3}{2},
\end{align*}
i.e.\ coefficient matrix
$\left(\begin{smallmatrix}7&0&11&5\\0&3&7&5\\2&0&1&1\end{smallmatrix}\right)$
(rows $=i=0,\ldots,2$; columns $=j=0,\ldots,3$), matching \texttt{balance} exactly.

\end{tcolorbox}

In many implementations the basis elements $L_{(s)}$ or ${\hat L}_{(s)}$ are stored as matrices.   For example, Eji\_LinComb
structures \url{https://github.com/kesterlester/Echinocoder/blob/main/Eji_LinComb.py} like those shown
below (only for each term without a zero coefficient) are used in
\url{https://github.com/kesterlester/Echinocoder/blob/main/C0HomDeg1_simplicialComplex_embedder_2a_for_array_of_reals_as_multiset.py}.
In these the count matrix has rows representing
coordinates $i{=}0,1,2$; and columns representing vectors $j{=}0,1,2,3$.
% AUTO-GENERATED by generate_simplex_example_tables.py — do not edit.
% Regenerate with: make tables   (or run the script directly).
\begin{center}
\small
\setlength{\tabcolsep}{4pt}
\renewcommand{\arraystretch}{1.2}
\begin{tabular}{cl l p{6.6cm}}
\toprule
$s$ & $\alpha_s$ & Eji\_LinComb structure
    &  ${\hat L}_{(s)}=L_{(s)}/(s{+}1)$ \\
\midrule
0 & 1
  & $\begin{pmatrix}1&0&1&1\\0&0&0&0\\0&0&0&0\end{pmatrix}$, idx 1
  & $\eji{0}{0}+\eji{2}{0}+\eji{3}{0}$ \\[8pt]
1 & 2
  & $\begin{pmatrix}1&0&2&1\\0&0&0&0\\0&0&0&0\end{pmatrix}$, idx 2
  & $\tfrac{1}{2} (\eji{0}{0}+2\eji{2}{0}+\eji{3}{0})$ \\[8pt]
2 & 3
  & $\begin{pmatrix}1&0&2&1\\0&1&1&1\\0&0&0&0\end{pmatrix}$, idx 3
  & $\tfrac{1}{3} (\eji{0}{0}+2\eji{2}{0}+\eji{3}{0}$\newline
    $\phantom{\tfrac{1}{3}(}$${}+\eji{1}{1}+\eji{2}{1}+\eji{3}{1})$ \\[8pt]
5 & 6
  & $\begin{pmatrix}2&0&3&1\\0&1&3&2\\0&0&0&0\end{pmatrix}$, idx 6
  & $\tfrac{1}{6} (2\eji{0}{0}+3\eji{2}{0}+\eji{3}{0}$\newline
    $\phantom{\tfrac{1}{6}(}$${}+\eji{1}{1}+3\eji{2}{1}+2\eji{3}{1})$ \\[8pt]
7 & 8
  & $\begin{pmatrix}2&0&3&1\\0&1&3&2\\2&0&1&1\end{pmatrix}$, idx 8
  & $\tfrac{1}{8} (2\eji{0}{0}+3\eji{2}{0}+\eji{3}{0}$\newline
    $\phantom{\tfrac{1}{8}(}$${}+\eji{1}{1}+3\eji{2}{1}+2\eji{3}{1}$\newline
    $\phantom{\tfrac{1}{8}(}$${}+2\eji{0}{2}+\eji{2}{2}+\eji{3}{2})$ \\[8pt]
\bottomrule
\end{tabular}
\end{center}


\noindent
(Terms with $\alpha_s=0$ contribute nothing and are omitted.)

\noindent
The table below repeats the construction with a different but equally valid
tie-breaking order: % AUTO-GENERATED by generate_simplex_example_tables.py — do not edit.
% Regenerate with: make tables   (or run the script directly).
within each group of equal gaps, the gap--basis-element assignment is rotated one place (for a two-element group, a swap).  The tied groups here are $\delta=5$ (at $s\in\{2,3\}$), $\delta=2$ (at $s\in\{5,6\}$) and $\delta=1$ (at $s\in\{7,8\}$).%

The basis vectors in the greyed cells change; the non-greyed cells
are identical to those in the table above.  Critically the greyed cells all have basis-coefficients
(scalars  $\Delta_{(s)}$ or $\alpha_{s}$) which are zero by construction ... and so
 the very places where the basis-vectors are ambiguous are also (by construction)
  the very places where those basis vectors do not need to be encoded. Equivalently the
   non-zero values of $\Delta_s$ or $\alpha_s$ all scale basis vectors that are unambiguously defined.
    This is the key mechanism by which the embedder is able to maintain continuity in its outputs despite
      having to use the outputs to both encode the coefficients (which change smoothly) and the basis vectors
      themselves (which change discontinuously at boundaries).

\input{generated/sx2_stepD_variant}

\subsection*{Step E: Canonicalise}

Sort the columns of each count matrix lexicographically (the column-sort
permutation ``forgets'' which vector was which, achieving multiset invariance).
After canonicalisation, $L_s$ becomes $\canon{L_s}$.

\subsection*{Step F: Hash and assemble}

Each $\canon{L_s}$ (along with its index) is hashed to a point $p_s \in
[0,1]^{19}$ (since $N_{\text{hash}} = 2{\times}9{+}1 = 19$).
The output is:
% AUTO-GENERATED by generate_simplex_example_tables.py — do not edit.
% Regenerate with: make tables   (or run the script directly).
\[
  f(M) \;=\;
  \bigl[\;\underbrace{-3,\;2,\;1}_{3\text{ row min.}},\;
          1\,p_{0} + 2\,p_{1} + 3\,p_{2} + 6\,p_{5} + 8\,p_{7}\;\bigr]
  \;\in\; \R^{22}.
\]


% ═════════════════════════════════════════════════════════════════════════════
\section{Worked Example: Algorithm 1, $n=4$, $k=2$}
\label{sec:example-alg1}
% ═════════════════════════════════════════════════════════════════════════════

\[
  M \;=\; \left\{\;
    \begin{pmatrix}4\\2\end{pmatrix},\;
    \begin{pmatrix}-3\\5\end{pmatrix},\;
    \begin{pmatrix}8\\9\end{pmatrix},\;
    \begin{pmatrix}2\\7\end{pmatrix}
  \;\right\}
  \quad (n=4 \text{ vectors in } \R^2)
\]

\subsection*{Step B: Extract global minimum (and maximum)}

Global min $= -3$, global max $= 9$.  Both stored as anchors.  Balance has
all non-negative entries.

\subsection*{Steps C--D: Global sort and Abel summation}

The balance after subtracting the global minimum $-3$ from every entry is:
\[
  \left\{\;\begin{pmatrix}7\\5\end{pmatrix},\;
            \begin{pmatrix}0\\8\end{pmatrix},\;
            \begin{pmatrix}11\\12\end{pmatrix},\;
            \begin{pmatrix}5\\10\end{pmatrix}\;\right\}
\]
Sorted globally descending (value: label):
\[
  12(\eji{2}{1}),\;11(\eji{2}{0}),\;10(\eji{3}{1}),\;8(\eji{1}{1}),\;
  7(\eji{0}{0}),\;5(\eji{0}{1}),\;5(\eji{3}{0}),\;0(\eji{1}{0}).
\]
The 7 gaps are $1,1,2,1,2,0,5$.  After global re-sort by gap descending and
second Abel summation: $\alpha_0=3,\;\alpha_2=3,\;\alpha_5=6$ (others zero).

$N_{\text{hash}} = 2(8{-}1){+}1 = 15$. Output:
\[
  f(M) \;=\;
  \bigl[\;3\,p_0 + 3\,p_2 + 6\,p_5,\;9,\;{-3}\;\bigr]
  \;\in\; \R^{17},
\]
numerically (from the code's debug output):
\[
  [\,7.07,\;6.91,\;2.78,\;3.50,\;1.01,\;5.51,\;4.07,\;4.90,\;10.02,\;
   9.51,\;5.13,\;2.97,\;8.04,\;5.34,\;7.41,\;9,\;{-3}\,].
\]

% ═════════════════════════════════════════════════════════════════════════════
\section{Output Size and the Whitney Embedding Theorem}
\label{sec:whitney}
% ═════════════════════════════════════════════════════════════════════════════

Both algorithms produce output of size $\approx 2nk$, which is not overhead to
be engineered away — it is intrinsic.  The input space $\SP{n}{k}$ is an
$nk$-dimensional manifold.  The \textbf{Whitney Embedding Theorem} states that
any smooth $d$-dimensional manifold embeds smoothly into $\R^{2d}$, and $2d$
is tight in general.  With output $\approx 2nk$, both algorithms sit right at
the Whitney bound.

The factor of $\approx 2$ is in fact the price of the cone-dimension mechanism
in Step~F: encoding $m$ positive coefficients together with the identities of
$m$ basis elements inside a single point requires $\R^{2m+1}$.  This same
mechanism is what delivers $C^0$ continuity (shared faces at zero-coefficient
transitions).  The factor of two is the minimum price of a smooth
embedding, and the algorithms pay almost exactly that.

Whether $2nk$ (or $2nk+1$, or $2nk+1-k$) is optimal for the specific
manifold $\SP{n}{k}$, or whether the known structure of that space admits a
tighter construction, remains open.
